\chapter{NỘI DUNG THỰC TẬP}

\textit{Lưu ý: Sinh viên cần cập nhật báo cáo theo từng tuần dựa trên công việc thực tế đã thực hiện.}

\section{Mô tả công việc thực tập}

Sinh viên mô tả tổng quan công việc được giao trong thời gian thực tập, bao gồm mục tiêu công việc, phạm vi công việc, công cụ sử dụng và sản phẩm/kết quả cần đạt được.

\section{Công việc thực tập}

Sinh viên trình bày chi tiết các nhóm công việc đã thực hiện. Ví dụ:

\begin{itemize}
    \item Tìm hiểu quy trình phát triển phần mềm tại doanh nghiệp.
    \item Cài đặt môi trường phát triển.
    \item Tìm hiểu tài liệu yêu cầu, tài liệu thiết kế hoặc mã nguồn.
    \item Tham gia xây dựng, kiểm thử hoặc bảo trì một chức năng cụ thể.
    \item Viết tài liệu hướng dẫn sử dụng hoặc tài liệu kỹ thuật.
\end{itemize}

\section{Chi tiết nhật ký thực tập}

Sinh viên có thể trình bày nhật ký thực tập theo bảng dưới đây.

\begin{longtable}{|c|p{3cm}|p{6cm}|p{4cm}|}
\caption{Nhật ký thực tập}
\label{tab:nhat-ky-thuc-tap} \\
\hline
\textbf{Tuần} & \textbf{Thời gian} & \textbf{Nội dung công việc} & \textbf{Kết quả đạt được} \\
\hline
\endfirsthead

\hline
\textbf{Tuần} & \textbf{Thời gian} & \textbf{Nội dung công việc} & \textbf{Kết quả đạt được} \\
\hline
\endhead

1 & Từ ngày ... đến ngày ...
& Tìm hiểu doanh nghiệp, phòng ban thực tập, quy trình làm việc.
& Nắm được cơ cấu tổ chức, nội quy và công việc được giao. \\
\hline

2 & Từ ngày ... đến ngày ...
& Cài đặt môi trường làm việc, tìm hiểu công nghệ sử dụng trong dự án.
& Cài đặt thành công môi trường, chạy thử dự án mẫu. \\
\hline

3 & Từ ngày ... đến ngày ...
& Tham gia thực hiện một số nhiệm vụ được giao trong dự án.
& Hoàn thành các công việc theo yêu cầu của người hướng dẫn. \\
\hline

4 & Từ ngày ... đến ngày ...
& Kiểm thử, chỉnh sửa lỗi, hoàn thiện tài liệu hoặc báo cáo kết quả.
& Hoàn thiện sản phẩm/công việc được giao. \\
\hline

\end{longtable}
\Nguon{Sinh viên tổng hợp}
